\documentclass[../DoAn.tex]{subfiles}
\begin{document}

Chương~1 đã xác định bài toán MAPF và đề xuất định hướng giải pháp. Chương này phân
tích chi tiết hiện trạng các giải pháp điều hướng đa agent trong game, từ đó xác định
yêu cầu chức năng và phi chức năng của hệ thống cần xây dựng.

% ============================================================
\section{Khảo sát hiện trạng}
\label{section:2.1}

\subsection{Các giải pháp điều hướng trong game hiện tại}

Các giải pháp điều hướng agent trong game thương mại và nghiên cứu học thuật có thể
phân thành bốn nhóm chính:

\textbf{(1) Unity NavMesh}: Unity Engine cung cấp hệ thống dẫn đường tự động dựa trên
navigation mesh 3D. NavMesh tính toán đường đi cho từng agent độc lập và hỗ trợ tránh
va chạm theo thời gian thực qua \textit{NavMeshAgent}. Ưu điểm: tích hợp sẵn, dễ sử
dụng, hỗ trợ môi trường 3D phức tạp. Nhược điểm: không giải quyết bài toán MAPF (không
đảm bảo tránh va chạm ở lớp lập kế hoạch cho môi trường lưới rời rạc), không tương
thích với benchmark \texttt{.map} chuẩn MAPF.

\textbf{(2) A* đơn agent}: Thuật toán A* \cite{hart1968astar} với heuristic Manhattan
trên lưới ô vuông là giải pháp phổ biến nhất trong game 2D. Mỗi agent tìm đường độc
lập mà không biết kế hoạch của agent khác. Ưu điểm: đơn giản, hiệu quả, dễ triển khai.
Nhược điểm: va chạm giữa các agent phải xử lý ở lớp vật lý, không có phối hợp ở lớp
lập kế hoạch, hiệu năng suy giảm khi số agent lớn do xử lý va chạm phản ứng.

\textbf{(3) CBS và thuật toán MAPF tối ưu}: Conflict-Based Search \cite{sharon2015cbs}
đảm bảo giải pháp tối ưu toàn cục (tổng số bước di chuyển nhỏ nhất). Ưu điểm: tối ưu,
hoàn chỉnh. Nhược điểm: độ phức tạp tính toán tăng theo hàm mũ với số agent và mật độ
va chạm, không thực tế cho game thời gian thực với nhiều agent.

\textbf{(4) PIBT}: Priority Inheritance with Backtracking \cite{okumura2022pibt} là thuật
toán MAPF chạy theo bước thời gian, hoàn chỉnh trên bản đồ có đường đi tồn tại, với
độ phức tạp thực tế phù hợp cho môi trường thời gian thực. PIBT không đảm bảo tối ưu
toàn cục nhưng thực hành cho kết quả tốt và đã được dùng trong các cuộc thi MAPF quốc
tế \cite{movingai2019benchmark}.

Bảng~\ref{tab:comparison} tổng hợp so sánh bốn giải pháp theo các tiêu chí phù hợp
với yêu cầu của đồ án.

\begin{table}[htbp]
\centering
\caption{So sánh các giải pháp điều hướng đa agent}
\label{tab:comparison}
\small
\begin{tabular}{lcccc}
\hline
\textbf{Tiêu chí}           & \textbf{NavMesh} & \textbf{A* đơn} & \textbf{CBS} & \textbf{PIBT} \\ \hline
Tránh va chạm đa agent     & Vật lý           & Vật lý          & Lập kế hoạch & Lập kế hoạch  \\
Tính tối ưu                & Không            & Không           & Có           & Không         \\
Thời gian thực             & Có               & Có              & Không        & Có            \\
Tương thích bản đồ .map    & Không            & Có              & Có           & Có            \\
Tích hợp Unity             & Có sẵn           & Tự triển khai   & Tự triển khai & Tự triển khai \\
Phù hợp đồ án              & Không            & Baseline        & Không        & Chính         \\ \hline
\end{tabular}
\end{table}

Từ phân tích trên, đồ án lựa chọn A* làm baseline để có điểm so sánh quen thuộc,
và PIBT làm thuật toán chính vì cân bằng tốt giữa hiệu quả và khả năng thực thi
thời gian thực.

\subsection{Khảo sát bản đồ benchmark MAPF}

MovingAI MAPF benchmark \cite{sturtevant2012benchmarks} cung cấp hàng nghìn bản đồ
với định dạng \texttt{.map} chuẩn, phân loại theo môi trường (mê cung, trong nhà, game).
Đồ án sử dụng 5 bản đồ đại diện cho 5 kịch bản điển hình (Bảng~\ref{tab:maps}).

\begin{table}[htbp]
\centering
\caption{Năm bản đồ benchmark sử dụng trong đồ án}
\label{tab:maps}
\small
\begin{tabular}{llrrl}
\hline
\textbf{Nhãn} & \textbf{Tệp bản đồ}            & \textbf{Kích thước} & \textbf{Ô đi được} & \textbf{Đặc điểm} \\ \hline
Alpha32   & random-32-32-10.map       & 32×32     & 922  (90\%)  & Thoáng, ít tường  \\
Mansion   & ht\_mansion\_n.map        & 133×270   & 8959 (25\%)  & Hành lang hẹp     \\
Chantry   & ht\_chantry.map           & 162×141   & 7461 (33\%)  & Phòng kín         \\
Gallows   & lt\_gallowstemplar\_n.map & 251×180   & 10021 (22\%) & Vùng hẹp nhiều     \\
Maze128   & maze-128-128-10.map       & 128×128   & 14818 (90\%) & Mê cung lớn       \\ \hline
\end{tabular}
\end{table}

% ============================================================
\section{Tổng quan chức năng}
\label{section:2.2}

\subsection{Biểu đồ use case tổng quan}
\label{subsection:2.2.1}

Hệ thống có hai tác nhân chính. Tác nhân thứ nhất là người chơi (Player), trực tiếp
điều khiển xe tăng, bắn đạn, chọn chế độ chơi, cũng như cấu hình và khởi động backtest.
Tác nhân thứ hai là agent AI (Enemy), được điều khiển bởi thuật toán MAPF để tự động tìm
đường và tấn công Eagle Base. Hai tác nhân này tương tác trong cùng một môi trường lưới,
nơi người chơi cố gắng bảo vệ Eagle Base còn các agent AI phối hợp tấn công.

Các chức năng chính của hệ thống gồm: (UC1)~Chơi đơn một người chống lại các agent AI,
(UC2)~chơi mạng LAN nhiều người cùng chống AI, (UC3)~chọn thuật toán AI (A*, PIBT,
PIBT TCP), (UC4)~chạy backtest tự động, và (UC5)~xem kết quả backtest. Trong đó UC1,
UC3 và UC4 được đặc tả chi tiết tại Mục~\ref{section:2.3}; UC2 và UC5 được đặc tả trong
Phụ lục~\ref{tab:uc2} và~\ref{tab:uc5}.

Hình~\ref{fig:activity_single} mô tả luồng hoạt động trong kịch bản chơi đơn,
từ khi tải bản đồ đến khi kết thúc ván chơi.

\begin{figure}[htbp]
    \centering
    \includegraphics[width=0.85\textwidth]{Hinhve/activity_single_play.pdf}
    \caption{Biểu đồ hoạt động kịch bản chơi đơn}
    \label{fig:activity_single}
\end{figure}

Hình~\ref{fig:activity_internet} mô tả luồng hoạt động trong kịch bản chơi mạng LAN,
bao gồm bước tạo/tham gia phòng và đồng bộ trạng thái giữa các máy.

\begin{figure}[htbp]
    \centering
    \includegraphics[width=0.85\textwidth]{Hinhve/activity_internet.pdf}
    \caption{Biểu đồ hoạt động kịch bản chơi mạng LAN}
    \label{fig:activity_internet}
\end{figure}

\subsection{Quy trình nghiệp vụ backtest}
\label{subsection:2.2.2}

Backtest là quy trình nghiệp vụ quan trọng nhất từ góc độ nghiên cứu. Luồng hoạt động:
người dùng cấu hình tham số (chọn bản đồ, thuật toán, số lần lặp, bật/tắt chướng ngại
động), hệ thống tự động lặp qua tất cả tổ hợp, mỗi run thu thập chỉ số thời gian thực,
sau đó xuất CSV và chart tổng hợp. Quy trình này được mô tả chi tiết trong Chương~4
tại Hình~\ref{fig:backtest_workflow}.

% ============================================================
\section{Đặc tả chức năng}
\label{section:2.3}

\subsection{Đặc tả UC1 -- Chơi đơn một người}

\begin{table}[htbp]
\centering
\caption{Đặc tả UC1 -- Chơi đơn}
\label{tab:uc1}
\small
\begin{tabular}{lp{10cm}}
\hline
\textbf{Tên use case}   & Chơi đơn (Single Play) \\
\textbf{Tác nhân}       & Người chơi \\
\textbf{Tiền điều kiện} & Người chơi đã mở ứng dụng và chọn chế độ chơi đơn \\
\textbf{Luồng chính}    &
1. Người chơi chọn bản đồ và thuật toán AI từ menu. \newline
2. Hệ thống đọc tệp \texttt{.map}, dựng bản đồ và spawn xe tăng người chơi + Eagle Base. \newline
3. \texttt{MapScenarioBootstrap} spawn các agent AI theo thuật toán đã chọn. \newline
4. Người chơi điều khiển xe tăng qua bàn phím/chuột; agent AI tự động di chuyển và tấn công Eagle. \newline
5. Ván chơi kết thúc khi Eagle bị phá hủy (agent thắng) hoặc người chơi tiêu diệt hết agent. \\
\textbf{Luồng phát sinh} &
2a. Tệp \texttt{.map} không tìm thấy → hệ thống dùng bản đồ mặc định Alpha32. \newline
5a. Không có agent sống, Eagle còn HP → người chơi thắng. \\
\textbf{Hậu điều kiện}  & Màn hình kết quả hiển thị điểm máu Eagle và thời gian ván chơi \\
\hline
\end{tabular}
\end{table}

\subsection{Đặc tả UC4 -- Chạy backtest tự động}

\begin{table}[htbp]
\centering
\caption{Đặc tả UC4 -- Chạy backtest tự động}
\label{tab:uc4}
\small
\begin{tabular}{lp{10cm}}
\hline
\textbf{Tên use case}   & Chạy backtest tự động \\
\textbf{Tác nhân}       & Người dùng (nghiên cứu viên) \\
\textbf{Tiền điều kiện} & Ứng dụng đang ở menu chính; nếu chọn PIBT\_TCP, server C++ đang chạy \\
\textbf{Luồng chính}    &
1. Người dùng mở UI Backtest, chọn bản đồ, thuật toán, số lần lặp, bật/tắt chướng ngại động. \newline
2. Người dùng bấm Start. \newline
3. \texttt{BacktestRunner} lặp qua từng tổ hợp (map × algorithm × rep). \newline
4. Mỗi run: spawn agent, chạy kịch bản đến khi Eagle bị phá hoặc timeout, ghi nhận chỉ số. \newline
5. Sau khi hết tất cả run: xuất \texttt{backtest\_summary\_*.csv}, \texttt{backtest\_agents\_*.csv}, \texttt{backtest\_chart\_*.html}. \\
\textbf{Luồng phát sinh} &
3a. PIBT\_TCP: nếu kết nối TCP thất bại, run đó ghi \texttt{Outcome=ConnectionError} và tiếp tục. \newline
4a. Run chạm timeout (mặc định 120\,s) → ghi \texttt{Outcome=Timeout}. \\
\textbf{Hậu điều kiện}  & Tệp CSV và HTML chart xuất hiện trong thư mục \texttt{BacktestResults/} \\
\hline
\end{tabular}
\end{table}

\subsection{Đặc tả UC3 -- Chọn thuật toán AI}

\begin{table}[htbp]
\centering
\caption{Đặc tả UC3 -- Chọn thuật toán AI}
\label{tab:uc3}
\small
\begin{tabular}{lp{10cm}}
\hline
\textbf{Tên use case}   & Chọn thuật toán AI \\
\textbf{Tác nhân}       & Người chơi hoặc nghiên cứu viên \\
\textbf{Tiền điều kiện} & Ứng dụng đang ở menu chọn chế độ chơi \\
\textbf{Luồng chính}    &
1. Người dùng chọn một trong ba thuật toán: A*, PIBT, PIBT\_TCP. \newline
2. Hệ thống ghi nhận lựa chọn vào \texttt{BacktestMode.Algorithm}. \newline
3. Khi bắt đầu ván: \texttt{MapScenarioBootstrap} spawn đúng loại agent tương ứng (\texttt{GridEnemyAgent}, \texttt{GridEnemyAgentPIBT}, hoặc \texttt{GridEnemyAgentPIBT\_TCP}). \newline
4. Nếu chọn PIBT\_TCP: hệ thống gửi lệnh \texttt{Hello} đến server, xác nhận kết nối trước khi spawn agent. \\
\textbf{Luồng phát sinh} &
4a. Server C++ không phản hồi trong 5\,s → báo lỗi và quay lại menu. \\
\textbf{Hậu điều kiện}  & Agent AI trong kịch bản dùng thuật toán đã chọn \\
\hline
\end{tabular}
\end{table}

% ============================================================
\section{Yêu cầu phi chức năng}
\label{section:2.4}

\textbf{Hiệu năng}: Thuật toán A* phải hoàn thành tìm đường cho mỗi agent trong dưới
5\,ms trên bản đồ nhỏ (Alpha32), để tổng thời gian tính toán mỗi khung hình không vượt
quá 16\,ms (60\,FPS). PIBT phải hoàn thành một bước phối hợp cho 6 agent trong dưới
10\,ms.

\textbf{Tính ổn định}: Mỗi run backtest phải tái lập được với cùng seed; hệ thống không
được crash trong suốt 90 run liên tiếp.

\textbf{Tính tương thích}: Ứng dụng chạy trên Windows 10/11 với Unity 6000.3.x; chế
độ WebGL tương thích Chrome và Firefox. Tệp \texttt{.map} phải đọc được theo chuẩn
MovingAI không cần chỉnh sửa.

\textbf{Khả năng mở rộng}: Thêm thuật toán mới chỉ cần tạo lớp agent mới kế thừa
interface chung, không cần thay đổi \texttt{BacktestRunner} hay \texttt{TankController}.

\textbf{Kết chương~2:} Chương này đã phân tích bốn nhóm giải pháp điều hướng đa agent
và xác nhận A* (baseline) cùng PIBT (thuật toán chính) là lựa chọn phù hợp nhất với yêu
cầu đồ án. Năm bản đồ benchmark chuẩn đã được chọn để đảm bảo tính so sánh được với
nghiên cứu cộng đồng. Các use case và yêu cầu phi chức năng được đặc tả là cơ sở cho
thiết kế hệ thống trong Chương~3 và Chương~4.

\end{document}
